\documentclass[12pt,a4paper]{article}
\usepackage[margin=2cm]{geometry}
\usepackage{xeCJK}
\usepackage{fontspec}
\setCJKmainfont{Noto Serif CJK TC}[Script=CJK]
\usepackage{amsmath,amssymb}
\usepackage{graphicx}
\usepackage{fancyhdr}
\setlength{\headheight}{14.5pt}
\addtolength{\topmargin}{-2.5pt}
\usepackage{hyperref}
\usepackage{listings}
\usepackage{enumitem}
\usepackage{titlesec}
\usepackage{caption}
\usepackage{indentfirst}
\usepackage{booktabs}
\usepackage{longtable}
\usepackage{multirow}
\usepackage{array}
\usepackage{tabularx}
\usepackage{float}
\usepackage{minted}
\usepackage{xcolor}
\setlength{\parindent}{2em}
\pagestyle{fancy}
\fancyhf{}
\cfoot{\thepage}
\linespread{1.3}
\setminted{
    linenos,                % 行號
    frame=lines,            % 上下框線
    framesep=5pt,           % 程式碼與邊框距離
    numbersep=8pt,          % 行號與程式碼距離
    fontsize=\scriptsize,   % 字體大小
    breaklines,             % 自動換行
    tabsize=4,              % tab 寬度
    rulecolor=\color{black},% 框線顏色
    xleftmargin=1.5em       % 左側縮排
}

\title{Machine Learning in Human Motion Analysis HW 1}
\author{B12508026戴偉璿}
\date{\today}

\begin{document}

\maketitle

\lhead{Machine Learning in Human Motion Analysis HW 1}
\rhead{B12508026戴偉璿}

\begin{enumerate}
    \item Find the extrema of the Matyas function using the optimality criteria method.
    
    {\color{blue}
    Solution:

    The Matyas function is defined as:
    $$f(x) = 0.26(x_1^2 + x_2^2) - 0.48x_1x_2$$

    To find the first order necessary condition, compute the gradient of the function:
    $$\nabla f(x) = \begin{bmatrix}0.52x_1 - 0.48x_2 \\ 0.52x_2 - 0.48x_1\end{bmatrix}$$

    Setting the gradient to zero gives the first-order necessary condition:
    \begin{align*}
    0.52x_1 - 0.48x_2 &= 0 \\
    0.52x_2 - 0.48x_1 &= 0
    \end{align*}

    Solving the above equations, we get:
    \begin{align*}
    x_1 &= 0 \\
    x_2 &= 0
    \end{align*}

    Check the second order sufficient condition by computing the Hessian matrix:
    $$H = \begin{bmatrix}0.52 & -0.48 \\ -0.48 & 0.52\end{bmatrix}$$

    Check the positive definiteness:
    \begin{align*}
    H_{11} & > 0 \\
    \text{det}(H) &= (0.52)(0.52) - (-0.48)(-0.48)  > 0
    \end{align*}

    Thus, $H$ is positive definite, and the second order sufficient condition is satisfied. Therefore, the point $(0, 0)$ is a local minimum of the Matyas function.
    }

    \item Find the extrema of the Himmelblau's function using the optimality criteria method.

    {\color{blue}
    Solution:

    The Himmelblau's function is defined as:
    $$f(x) = (x_1^2 + x_2 - 11)^2 + (x_1 + x_2^2 - 7)^2$$
    
    To find the first order necessary condition, compute the gradient of the function:
    \begin{align*}
    \nabla f(x) &= \begin{bmatrix}4x_1(x_1^2 + x_2 - 11) + 2(x_1 + x_2^2 - 7) \\ 2(x_1^2 + x_2 - 11) + 4x_2(x_1 + x_2^2 - 7)\end{bmatrix}
    \end{align*}

    Setting the gradient to zero gives the first-order necessary condition:
    \begin{align*}
    4(x_1^2 + x_2 - 11)x_1 + 2(x_1 + x_2^2 - 7) &= 0 \\
    2(x_1^2 + x_2 - 11) + 4(x_1 + x_2^2 - 7)x_2 &= 0
    \end{align*}

    Solving the above equations, we find the following critical points:
    \begin{align*}
    (3.0, 2.0), \quad (-2.805, 3.131), \quad (-3.779, -3.283), \quad (3.584, -1.848), \quad (−0.271,−0.923)\\
    (-0.128,-1.954),\quad (0.087,2.884),\quad (3.385,0.0739),\quad (-3.073,-0.081).
    \end{align*}
    
    Check the second order sufficient condition by computing the Hessian matrix:
    \begin{align*}
    H = \begin{bmatrix}12x_1^2 + 4x_2 - 42 & 4x_1 + 4x_2 \\ 4x_1 + 4x_2 & 4x_1 + 12x_2^2 - 26\end{bmatrix}
    \end{align*}
    \begin{align*}
    \text{det}(H) =& (12x_1^2 + 4x_2 - 42)(4x_1 + 12x_2^2 - 26) - (4x_1 + 4x_2)^2\\
        =& 48x_1^3 + 144x_1^2x_2^2 - 328x_1^2 -16x_1x_2 + 48x_2^3 - 104x_2 - 168x_1 - 520x_2^2 + 1092
    \end{align*}


    Apply each critical point into the Hessian matrix to check the second order sufficient condition:

    For points $(3.0, 2.0)$, $(-2.805, 3.131)$, $(-3.779, -3.283)$, and $(3.584, -1.848)$, their $H_{11}$ values are positive and their determinants are positive, thus they are local minima of the Himmelblau's function.

    For $(-0.271, -0.923)$, its $H_{11}$ value is negative and its determinant is positive, thus it is a local maximum of the Himmelblau's function.

    For points $(-0.128,-1.954)$, $(0.087,2.884)$, $(3.385,0.0739)$, and $(-3.073,-0.081)$, their $H_{11}$ values are negative and their determinants are negative, thus they are saddle points of the Himmelblau's function.


    }
\end{enumerate}

\end{document}